\subsection{\problemblock{*Atmos} Data Block}\label{sec:block-atmos}

An atmosphere model is required to integrate the equations of motion. It defines
temperature, pressure, density, speed of sound, and viscosity as functions of altitude,
as described in \cref{sec:atmosphere-model}. TAOS contains the 20 standard
models listed in \cref{tab:atmosphere-types} and also supports two forms of
user-defined atmosphere. The default is the 1976 U.S. Standard Atmosphere.

The models extend to 1000~km. Above 146~km, all models use the same values as the
1976 U.S. Standard Atmosphere. Below sea level, temperature is extrapolated using
the sea-level temperature gradient and pressure is held at the sea-level value; the
other properties are computed from pressure and temperature.

A standard atmosphere is selected by placing its number after the
\problemblock{*atmos} keyword. For example,

\begin{lstlisting}[style=taosmanual]
*atmos 20
\end{lstlisting}

selects the Kwajalein mean-annual atmosphere. One selected atmosphere applies to
all trajectories in the problem; different trajectories cannot use different atmosphere
models. Because type 1 is common, the keyword \texttt{standard} may be used in
place of the number 1. Similarly, \texttt{none} may be used in place of type 0.

\begin{table}[htbp]
  \centering
  \caption{Atmosphere types.}
  \label{tab:atmosphere-types}
  \small
  \begin{tabularx}{0.94\linewidth}{@{}rXrX@{}}
    \toprule
    Type & Atmosphere & Type & Atmosphere \\
    \midrule
     0 & None & 11 & $60^\circ$ North July \\
     1 & 1976 U.S. Standard & 12 & $75^\circ$ North January \\
     2 & $15^\circ$ North Annual & 13 & $75^\circ$ North January (cold) \\
     3 & $30^\circ$ North January & 14 & $75^\circ$ North January (warm) \\
     4 & $30^\circ$ North July & 15 & $75^\circ$ North July \\
     5 & $45^\circ$ North January & 16 & Tonopah winter \\
     6 & $45^\circ$ North July & 17 & Tonopah spring \\
     7 & $45^\circ$ North spring/fall & 18 & Tonopah summer \\
     8 & $60^\circ$ North January & 19 & Tonopah fall \\
     9 & $60^\circ$ North January (cold) & 20 & Kwajalein mean annual \\
    10 & $60^\circ$ North January (warm) & & \\
    \bottomrule
  \end{tabularx}
\end{table}

\taossubhead{User-Defined Atmosphere}

The first user-defined form requires tabulated values for every atmospheric property
as a function of altitude. The format is illustrated by

\begin{lstlisting}[style=taosmanual]
*atmos user
  alt    temp   pres    rho       sndspd  visc
      0  548.2  2111.0  0.00224   1147   3.895e-7
   5000  530.0  1774.0  0.00195   1128   3.780e-7
  10000  511.0  1480.0  0.00168   1109   3.701e-7
  15000  494.9  1229.0  0.00124   1072   2.602e-7
  20000  477.5   833.0  0.00105   1051   3.400e-7
\end{lstlisting}

The keyword \texttt{user} indicates a user-defined atmosphere with all properties specified. The next
line contains column headers for altitude and all atmospheric properties. Every
keyword is required, but the order is arbitrary. These keywords are the same as the
atmospheric output variables in the Appendix, and their units may be changed with
\problemblock{*units/fmt}.

Atmospheric properties are tabulated against altitude, so altitude is normally the
first column. All properties on one line correspond to that altitude. Altitude values
must be strictly increasing; duplicate values are not allowed. During integration the
properties are linearly interpolated. Enough points must be provided for linear
interpolation to represent the desired atmosphere adequately. There are no limits on the number of values that can be input.

\taossubhead{Site-Measured Atmosphere}

The second user-defined form requires only pressure and density as functions of
altitude. It is called a site-measured atmosphere, and exponential interpolation is
used so fewer points can represent a realistic profile.

\begin{lstlisting}[style=taosmanual]
*atmos site
  alt    pres    rho
      0  2115.0  0.00223
  10000  1485.0  0.00166
  20000   830.0  0.00105
  30000   627.0  0.00089
  40000   241.0  0.00058
\end{lstlisting}

A site-measured atmosphere is input in the same way as a complete user-defined
atmosphere, except that there are fewer columns of data. Temperature is estimated
from pressure and density with the ideal-gas law, and speed of sound and viscosity
are computed from that temperature.
